\documentclass[12pt]{article}
\usepackage{latexblog}

% ============================================================
% Blog Metadata
% ============================================================
\blogtitle{使用GOCD部署一个SpringBoot应用}
\blogdate{2017-06-26}
\blogtags{Docker, 云计算, 容器化}
\bloglang{zh}
\blogsummary{}

\begin{document}
\makeblogtitle

如何自动部署 springboot 的应用？

首先我们要做的是把GOCD的启动用户\texttt{go}加入到\texttt{docker}组中，这样就可以免\texttt{sudo}来执行docker的命令了

\begin{Shaded}
\begin{Highlighting}[]
\FunctionTok{sudo}\NormalTok{ gpasswd }\AttributeTok{{-}a}\NormalTok{ go docker}
\FunctionTok{sudo}\NormalTok{ /etc/init.d/docker restart}
\ExtensionTok{newgrp} \AttributeTok{{-}}\NormalTok{ docker}
\end{Highlighting}
\end{Shaded}

我们现在来通过GOCD部署一个eureka的服务端。eureka用作服务发现，包含server端和client端，每个微服务是一个client注册到server上。看看需要哪些包：

\begin{Shaded}
\begin{Highlighting}[]
\NormalTok{dependencies }\OperatorTok{\{}
\NormalTok{    testCompile group}\OperatorTok{:} \StringTok{\textquotesingle{}junit\textquotesingle{}}\OperatorTok{,}\NormalTok{ name}\OperatorTok{:} \StringTok{\textquotesingle{}junit\textquotesingle{}}\OperatorTok{,}\NormalTok{ version}\OperatorTok{:} \StringTok{\textquotesingle{}4.12\textquotesingle{}}
\NormalTok{    compile group}\OperatorTok{:} \StringTok{\textquotesingle{}org.springframework.boot\textquotesingle{}}\OperatorTok{,}\NormalTok{ name}\OperatorTok{:} \StringTok{\textquotesingle{}spring{-}boot{-}starter{-}web\textquotesingle{}}\OperatorTok{,}\NormalTok{ version}\OperatorTok{:} \StringTok{\textquotesingle{}1.5.3.RELEASE\textquotesingle{}}
\NormalTok{    compile group}\OperatorTok{:} \StringTok{\textquotesingle{}org.springframework.cloud\textquotesingle{}}\OperatorTok{,}\NormalTok{ name}\OperatorTok{:} \StringTok{\textquotesingle{}spring{-}cloud{-}starter{-}eureka{-}server\textquotesingle{}}\OperatorTok{,}\NormalTok{ version}\OperatorTok{:} \StringTok{\textquotesingle{}1.3.1.RELEASE\textquotesingle{}}
\OperatorTok{\}}
\end{Highlighting}
\end{Shaded}

然后我们的程序写一句话就可以了：

\begin{Shaded}
\begin{Highlighting}[]
\AttributeTok{@EnableEurekaServer}
\AttributeTok{@SpringBootApplication}
\KeywordTok{public} \KeywordTok{class}\NormalTok{ EurekaServerApplication }\OperatorTok{\{}
    \KeywordTok{public} \DataTypeTok{static} \DataTypeTok{void} \FunctionTok{main}\OperatorTok{(}\BuiltInTok{String}\OperatorTok{[]}\NormalTok{ args}\OperatorTok{)\{}
\NormalTok{        SpringApplication}\OperatorTok{.}\FunctionTok{run}\OperatorTok{(}\NormalTok{EurekaServerApplication}\OperatorTok{.}\FunctionTok{class}\OperatorTok{,}\NormalTok{ args}\OperatorTok{);}
    \OperatorTok{\}}
\OperatorTok{\}}
\end{Highlighting}
\end{Shaded}

这样启动后就可以访问到eureka的web页面了，我们可以在application.properties中配置一些属性：

\begin{Shaded}
\begin{Highlighting}[]
\NormalTok{server.port=8761}
\NormalTok{eureka.client.register{-}with{-}eureka=false}
\NormalTok{eureka.client.fetch{-}registry=false}
\end{Highlighting}
\end{Shaded}

然后来想办法把应用打包到docker中。我们新建一个Dockerfile到src/main/docker中：

\begin{Shaded}
\begin{Highlighting}[]
\ConstantTok{FROM} \VariableTok{oracle}\OperatorTok{/}\VariableTok{openjdk}\OperatorTok{:}\DecValTok{8}
\ConstantTok{VOLUME} \OperatorTok{/}\VariableTok{tmp}
\ConstantTok{ADD} \VariableTok{eureka}\OperatorTok{{-}}\VariableTok{server}\OperatorTok{{-}}\DecValTok{1.0}\OperatorTok{{-}}\ConstantTok{SNAPSHOT}\OperatorTok{.}\VariableTok{jar} \VariableTok{app}\OperatorTok{.}\VariableTok{jar}
\ConstantTok{RUN} \VariableTok{sh} \OperatorTok{{-}}\NormalTok{c }\StringTok{\textquotesingle{}touch /app.jar\textquotesingle{}}
\ConstantTok{ENV} \ConstantTok{JAVA\_OPTS}\OperatorTok{=}\StringTok{""}
\ConstantTok{ENTRYPOINT} \OperatorTok{[} \StringTok{"sh"}\OperatorTok{,} \StringTok{"{-}c"}\OperatorTok{,} \StringTok{"java $JAVA\_OPTS {-}Djava.security.egd=file:/dev/./urandom {-}jar /app.jar"} \OperatorTok{]}
\end{Highlighting}
\end{Shaded}

然后在build.gradle中添加一个构建任务：

\begin{Shaded}
\begin{Highlighting}[]
\NormalTok{buildscript }\OperatorTok{\{}
\NormalTok{    repositories }\OperatorTok{\{}
        \FunctionTok{mavenCentral}\OperatorTok{()}
    \OperatorTok{\}}
\NormalTok{    dependencies }\OperatorTok{\{}
        \FunctionTok{classpath}\OperatorTok{(}\StringTok{"org.springframework.boot:spring{-}boot{-}gradle{-}plugin:1.5.2.RELEASE"}\OperatorTok{)}
        \FunctionTok{classpath}\OperatorTok{(}\StringTok{\textquotesingle{}se.transmode.gradle:gradle{-}docker:1.2\textquotesingle{}}\OperatorTok{)}
    \OperatorTok{\}}
\OperatorTok{\}}
\NormalTok{task }\FunctionTok{buildDocker}\OperatorTok{(}\NormalTok{type}\OperatorTok{:}\NormalTok{ Docker}\OperatorTok{,}\NormalTok{ dependsOn}\OperatorTok{:}\NormalTok{ build}\OperatorTok{)} \OperatorTok{\{}
\NormalTok{    push }\OperatorTok{=} \KeywordTok{false}
\NormalTok{    applicationName }\OperatorTok{=}\NormalTok{ jar}\OperatorTok{.}\NormalTok{baseName}
\NormalTok{    dockerfile }\OperatorTok{=} \FunctionTok{file}\OperatorTok{(}\StringTok{\textquotesingle{}src/main/docker/Dockerfile\textquotesingle{}}\OperatorTok{)}
\NormalTok{    doFirst }\OperatorTok{\{}
\NormalTok{        copy }\OperatorTok{\{}
\NormalTok{            from jar}
\NormalTok{            into stageDir}
        \OperatorTok{\}}
    \OperatorTok{\}}
\OperatorTok{\}}
\end{Highlighting}
\end{Shaded}

这样当我们执行\texttt{gradle\ build\ buildDocker}命令的时候就会把springboot应用打包成docker镜像了。


\end{document}
